% ============================================================
% Chapter 8: Sustainability and Ethical Implications
%   Section: Ethics of Data Management and Analysis
%     Subsection: Personal Data and Privacy
% Included from chapters/08-sustainability-ethics/02-data-ethics/data-ethics.tex
% ============================================================

\subsection{Personal Data and Privacy}

The system processes no data about its users: there are no accounts,
no profiles and no tracking, a consequence of cutting them from the
scope. The personal data it does process are the names of people that
appear in published news. The entity recogniser extracts them, and the
graph stores them as entity nodes linked to the articles and claims
that mention them, beside the claimant of each labelled fact.

These are data that a publisher has already made public, processed
for journalistic and research purposes, for which the General Data
Protection Regulation provides room for exemptions and derogations
(Article~85 for freedom of expression and information, Article~89 for
scientific research, each subject to safeguards); they are not combined
with any other source about the person, and nothing in the system
profiles individuals. A public deployment would still need its own
assessment of the legal basis. Three precautions follow nonetheless. The
graph and the lake are local and not published; the graph's query
console is read-only and behind an API key; and an entity is keyed by
its type and name only, with no attempt to identify the person beyond
what the article says. A claim attributed to a named person carries a
verdict about the \emph{claim}, not about the person, and the reader
view presents it that way.
